
\chapter{Computer Architecture}
\label{chap:computer-architecture}

\chapterintro{Computer architecture explains how the main parts of a computer cooperate to execute programs.}
\learningobjectives{\begin{itemize}\item Identify the CPU, memory, storage, and buses.\item Explain the fetch-decode-execute cycle.\item Understand why GPUs matter for AI.\end{itemize}}

\section{The Big Picture}
A computer is not a single object. It is a collection of cooperating components. The CPU executes instructions. Memory stores temporary working data. Storage keeps long-term files. Input and output devices communicate with the outside world.

\section{The CPU}
The Central Processing Unit is the general-purpose instruction executor of the computer.
\begin{definition}The CPU is the component that fetches, decodes, and executes program instructions.\end{definition}

\section{The Fetch-Decode-Execute Cycle}
\begin{enumerate}\item Fetch the next instruction from memory.\item Decode what the instruction means.\item Execute the instruction.\item Store or communicate the result.\end{enumerate}

\section{Why GPUs Matter}
A GPU contains many smaller processing units designed for parallel numerical work. Deep learning uses enormous matrix operations, so GPUs are often much faster than CPUs for neural networks.

\section{Exercise}
\begin{exercise}Explain why a GPU can be useful for training neural networks.\end{exercise}
\begin{solution}Neural networks require many similar numerical operations on matrices and tensors. GPUs are designed to perform many operations in parallel, making them well suited to this workload.\end{solution}

\section{Summary}
Computer architecture connects software instructions to physical hardware. LLM engineering depends on this because model training and inference are constrained by processors, memory, storage, and networking.
